\documentclass[11pt,a4paper]{article}
\usepackage[margin=26mm]{geometry}
\usepackage{amsmath,amssymb,amsthm,mathtools}
\usepackage{lmodern,microtype,booktabs,longtable,array}
\usepackage[colorlinks=true,linkcolor=blue,citecolor=blue,urlcolor=blue]{hyperref}
\hypersetup{pdftitle={A small improvement of the seven-eighths zero-free boundary from a compensated probe},pdfauthor={RH-Weil research project}}
\newtheorem{theorem}{Theorem}[section]
\newtheorem{proposition}[theorem]{Proposition}
\newtheorem{lemma}[theorem]{Lemma}
\newtheorem{corollary}[theorem]{Corollary}
\theoremstyle{definition}
\newtheorem{definition}[theorem]{Definition}
\newtheorem{remark}[theorem]{Remark}
\newcommand{\Rea}{\operatorname{Re}}
\newcommand{\Ima}{\operatorname{Im}}
\newcommand{\N}{\operatorname{N}}
\newcommand{\eps}{\varepsilon}
\newcommand{\sig}{\sigma_1}
\newcommand{\bet}{\beta_*}
\newcommand{\HH}{\mathfrak H}
\newcommand{\LL}{L_F^S}
\newcommand{\pos}[1]{\left(#1\right)_+}
\newcommand{\src}[1]{\nolinkurl{#1}}
\makeatletter
\g@addto@macro\UrlBreaks{\do\a\do\b\do\c\do\d\do\e\do\f\do\g\do\h\do\i\do\j\do\k\do\l\do\m\do\n\do\o\do\p\do\q\do\r\do\s\do\t\do\u\do\v\do\w\do\x\do\y\do\z\do\0\do\1\do\2\do\3\do\4\do\5\do\6\do\7\do\8\do\9}
\makeatother
\title{A small improvement of the seven-eighths zero-free boundary\\from a compensated probe}
\author{RH--Weil research project}
\date{October 7, 2026\\[3pt]\small A proof relative to explicitly cited arithmetic and analytic inputs}
\begin{document}
\maketitle
\begin{abstract}
We give a parameter extension of the compensated-probe argument in the
September 30, 2026 version of \emph{The Quasi-Riemann Hypothesis}.
Relative to its specified generic arithmetic and analytic results and its
finite-order Hecke seven-eighths theorem, we obtain the strict zero-free
half-plane $\Rea s>69999/80000$ for finite-order Hecke $L$-functions over
$\mathbb Q(\sqrt{-3})$, and hence for all Dirichlet $L$-functions.
The total prime-slot length is increased from $1/6$ to
$1/6+1/20000$. A retained positive loss in the reflected row energy is
paid by the rescaling coefficients. We check the enlarged Euler domain,
the actual detector capacities, and the joint conductor cost of all
ramified errors. An exact continuous inequality gives a uniform central
saving $307/11016000$. We then specify the order of all choices, close
the height truncation, and use Mellin continuation for the full character
family. The assertion permits the principal pole at $s=1$ and does not
include the boundary line. This paper is a cited-input derivation;
it is not an independent certification of the source or a Lean
formalization of the improvement.
\end{abstract}
\tableofcontents

\section{Statement, provenance, and imported results}
\label{sec:inputs}
Put $F=\mathbb Q(\sqrt{-3})$ and $q_a=\N(a)$, with the source convention
for ideals, elements, ray representatives, and units. Characters are
extended by zero on nonunits. The superscript $S$ deletes the finite
Euler factors at the fixed set $S$. All assertions about a principal
$L$-function allow its pole at $1$.

Our source is the September 30 preprint \cite{OAI}, at commit
\src{adc7f1241b42e322a6451854ab7e4b4c146bf78a}. Its canonical UTF-8
SHA-256, after replacing CRLF and lone CR by LF, is recorded in
Appendix~\ref{app:verification}. The line references below refer to that
exact source, rather than to a moving branch.

\begin{definition}[Imported input package $\mathcal R$]
\label{def:R}
The package $\mathcal R$ consists of the following results with their
stated uniformities and hypotheses in \cite{OAI}: the coefficientwise
probe and Poisson identities and local arithmetic tables; smooth
separation, Gaussian annuli, reflection sectors and generic reflected
energy; the generic additive Gram bound; the global and buffered
reciprocal estimates, Hecke growth and functional equation; the
detector witnesses, marked and plain moments, sixth-power inverse
amplification, and uniform slot mesh; the fixed-ray prime asymptotic;
and the source theorem for all primitive finite-order Hecke characters
with strict boundary $7/8$.
The locations and the uses of these inputs are given in
Table~\ref{tab:inputs}. Their correctness is assumed here.
\end{definition}

\begin{theorem}[Boundary improvement relative to $\mathcal R$]
\label{thm:main}
Assume $\mathcal R$. For every finite-order Hecke character $\eta$ of
$F$, including imprimitive characters,
\[
 L_F(s,\eta)\ne0\qquad\left(\Rea s>\frac{69999}{80000}\right).
\]
For every Dirichlet character $\chi$ the same assertion holds for
$L(s,\chi)$, and in particular for $\zeta(s)$. The principal pole
at $s=1$ is allowed. No assertion is made about zeros on the line
$\Rea s=69999/80000$.
\end{theorem}

This is a change of $1/80000$ in the boundary of the cited theorem.
The proof does not use a proportion of zeros on the critical line.
It makes no claim about the Riemann hypothesis or the novelty of the
generic inputs.

\begingroup\small
\begin{longtable}{@{}>{\raggedright\arraybackslash}p{.22\textwidth}>{\raggedright\arraybackslash}p{.25\textwidth}>{\raggedright\arraybackslash}p{.45\textwidth}@{}}
\caption{Exact imported inputs and the new checks performed here.
Source labels and line intervals refer to \cite{OAI}.}\label{tab:inputs}\\
\toprule
Input & Source location & Use and retained hypotheses\\
\midrule
\endfirsthead
\toprule Input & Source location & Use and retained hypotheses\\ \midrule
\endhead
\bottomrule\endfoot
Smooth calculus and Gaussian annuli &
\src{lem:smooth-calculus}, \src{lem:gaussian-annular};
1123--1209, 1288--1315 &
Fixed finite-dimensional common Fourier profile, weighted seminorms,
and whole annular tails; no restriction to slot length $1/6$.\\
Reflection and energy &
\src{lem:reflection-row-sectors}, \src{lem:reflected-energy};
2470--2506, 7663--7804 &
Actual row masks, powerful fixed parts, retained dual support and
bounded logarithmic lengths. Section~\ref{sec:low} pays the changed loss.\\
Additive Gram &
\src{prop:probe-gram}; 8340--8364 &
Actual ray/Gauss coefficients, fixed polynomial ranges $Q,Y'\ge1$,
and $P_a=Y'^2/Q\ge1$.\\
Physical high identity &
\src{def:stage-high-data}; 5582--5649;
local compensation 8707--9049 &
An independently defined completed probe with finite compensation.
Section~\ref{sec:euler}
proves its new-domain correction without dividing by selected factors.\\
Reciprocal and contours &
400--502, 5672--5728, 5808--5896, 6021--6155 &
Whole buffered bins, all-height outer tails, and fixed real boxes.
No contour crosses a target zero.\\
Principal and outer rows &
6174--6301, 6342--6463, 9101--9195 &
Same finite $S$, same principal normalizer; independent absolute bounds
for small and large rows.\\
Detectors and moments &
\src{lem:detector-dyads}; 4385--4547;
marked 9220--9269; inverse amplification 12362--12389;
plain 12492--12578 &
Actual witnesses, both common signs, allowed heights and capacities,
and $\kappa\ge3/4$. Section~\ref{sec:count} checks each restriction.\\
Actual slot counts &
15015--15446 &
Whole slots, natural coefficients, inducing-character exclusion,
fixed finite presentation set, rounding, and common mesh.\\
Bootstrap and transfer &
16455--16463; 6705--6780 &
Whole Hecke $7/8$ theorem used only to bound $\bet$ from above;
quadratic transfer applied after the improved Hecke conclusion.\\
Height and order &
6520--6580, 16191--16453 &
Internal derivative orders fixed before the external tail order;
target-dependent height exponent chosen last.\\
\end{longtable}
\endgroup

\section{The same compensated physical expression}
\label{sec:physical}
The arithmetic data $S,T,b_*,\xi$, completed-index convention $A=cn^3$,
Gaussian test, ray calibration, and all zero-on-nonunit masks are
those of the source. Write $I_\eta(X,Y,Z)$ for its completed probe
(\src{eq:general-probe}), and $I_\eta^{[p\mid A]}$ for the same
expression restricted to $p\mid A$. The completed sums are independently
defined and absolutely convergent on the initial Mellin line; their
outer annular sums and the compensation over prime tuples and subsets
are finite. They are not defined by a high integral.

For a fixed prime tuple the two operations in
\src{eq:compensated-probe-definition} are the marked operation
$\overline{\eta(p)}I_\eta^{[p\mid A]}(X,Y,Zq_p)$ and the rescaled
operation $q_p^{-3/2}I_\eta(X/q_p,Y/q_p,Z)$, with a minus sign
between them. Apply these operations in each of the disjoint slots,
multiply by $W_i(q_p/P_i)$, and sum over the tuples. Denote the resulting
finite combination by $I_{\eta,\mathrm{mod}}(Z)$. The inherited
restrictions and calibration remain part of this definition.
More explicitly, if $p_J=\prod_{i\in J}p_i$, $p_\varnothing=1$, and
$I_{\eta;D}$ inserts $1_{D\mid cn^3}$ in the completed row, this is
\begin{equation}\label{eq:physicalsum}
\begin{split}
 I_{\eta,\mathrm{mod}}(Z)
 =\sum_{(p_i)}\prod_i W_i(q_{p_i}/P_i)
 \sum_{J\subseteq\{1,\ldots,K\}}&
 (-1)^{|J|}q_{p_J}^{-3/2}\overline{\eta(p_{J^c})}\\
 &\cdot I_{\eta;p_{J^c}}(X/q_{p_J},Y/q_{p_J},Zq_{p_{J^c}}).
\end{split}
\end{equation}
Every window stays at its original scale $P_i$, including when
another slot changes the completed scale.

We use the following fixed geometry:
\begin{equation}\label{eq:geometry}
\begin{gathered}
 b=\frac18,\quad \ell=\frac{10003}{60000}=\frac16+\frac1{20000},
 \quad l_x=\frac{1-\ell-b}{2}=\frac{42497}{120000},\\
 l_y=\frac{1-\ell+b}{2}=\frac{57497}{120000},\quad
 M=1-\ell,\quad h=1-l_x+\ell=\frac{32503}{40000},\\
 X=Z^{l_x},\quad Y=Z^{l_y},\quad
 \sig=\frac{69999}{80000},\qquad
 C(s)=s+\frac{l_x}{2}-1+\frac h6=s-\frac{11}{16}.
\end{gathered}
\end{equation}
In particular $C(\sig)=14999/80000$.
Take a fixed even integer $K$, positive equal lengths $\ell_i=\ell/K$,
and $P_i=Z^{\ell_i}$. Choose nonnegative nonzero smooth $W_i$ with
compact support in the pairwise disjoint intervals
\[
 \left(1+\frac{2i-1}{2K+1},\,1+\frac{2i}{2K+1}\right)\subset(1,2).
\]
The slots are the allowed source prime representatives in $T$ with
$q_p/P_i$ in these supports. Their disjointness makes every selected
prime occur once.

\section{The low bound at variable slot length}
\label{sec:low}
\begin{proposition}\label{prop:low}
Under the generic low inputs in $\mathcal R$, for every fixed
$1/6\le\ell\le1/5$, $b=1/8$, and the corresponding geometry above,
\[
 |I_{\eta,\mathrm{mod}}(Z)|
 \ll_{\eta,\mathrm{fixed},\eps}
 Z^{(1-\ell)/4-b/6+\eps}.
\]
The exponent loss $\eps>0$ can be prescribed before the target;
constants and thresholds may depend on its fixed arithmetic data.
\end{proposition}
\begin{proof}
Fix a rescaled subset $J$ and its prime tuple, and put
$d_J=\sum_{i\in J}\ell_i$, $\ell'=\ell-d_J$, $M'=M-2d_J$.
The norm ratio $r_J=q_{p_J}/Z^{d_J}$ lies in a fixed compact
subinterval of $(0,\infty)$. The exact low separation has a common
row ball
\[
 X'=XZ^{-d_J}/r_J,\quad Y'=YZ^{-d_J}/r_J,\quad
 Q=q_{b_*}X'Y'\asymp Z^{M'}.
\]
All surviving marked slots are summed inside $B_m^J$ before the
row norm is taken. The completed index remains $cn^3$, with weight
$q_c^{-1/2}q_n^{-1}$, including permitted shared primes.

Partition the Gaussian in $q_{cn^3}/Z^{1+\ell'}$ into whole annuli.
On $\log(q_{cn^3}/Z^{1+\ell'})=k+O(1)$ the common profile is
\[
 \chi(\log x)\,
 W_G\!\left(e^kx/\prod_i\varrho_i\right)\prod_i W_i(\varrho_i).
\]
Its weighted Fourier seminorms are summable over $k$. First use
the absolute completed-sum bound and Gaussian decay to remove
$|k|>\kappa_G\log Z+O(1)$ with arbitrary power saving.
The retained norm exponent is
$N_*=1+\ell'+\theta_N$,
$|\theta_N|\le\kappa_G+O(1/\log Z)$.
Fix this common Fourier density before separating rows or tuples.
The separated slot coefficients retain their arithmetic phases,
annular cutoffs and unit-modulus norm powers, so the generic
reflected-energy lemma applies.

Freeze a powerful row part of logarithmic length $O$ and an actual
residual squarefree dyad of length $H$. The source sectors give
\[
 \Delta_H=M'-O-H\ge-\eps_Z,\quad
 2A_0\le O+\eps_Z,\quad N_0\le A_0,\quad z_a\le\ell',
 \qquad \eps_Z=O_{\mathrm{fixed}}(1/\log Z).
\]
The dual length is
$T_d=2H+2A_0+2z_a-1-\ell'-\theta_N-N_0-3B_0$.
Since $M'+\ell'-1=-3d_J$,
\begin{equation}\label{eq:dual}
 T_d-(H-3d_J)
 =H-M'+2A_0+2(z_a-\ell')-N_0-3B_0-\theta_N
 \le2\eps_Z+|\theta_N|.
\end{equation}
This calculation retains the original support and shared $S$-primes.
For completeness the imported energy expression, after counting
the frozen parts once, is
\[
\begin{split}
 s_{\mathrm{hyb}}&=\min\{v,z_a,(v+z_a)/3\},\qquad
 y=v+3\ell_b+e_\lambda\le T_d+\tau_{\mathrm{ref}},\\
 E_{\mathrm{ref}}&=O/2+\max(H,v+\ell_b)-S_0-B_0+z_a
 -s_{\mathrm{hyb}}-\ell_b-2e_\lambda/3-(T_d-y)_+/2 .
\end{split}
\]
Here $v,\ell_b$ are the nonnegative dual ideal lengths,
$S_0,B_0,N_0,A_0$ are the frozen-part lengths of the source sectors,
and a negative $e_\lambda$ is only $O(1/\log Z)$.
The retained support and \eqref{eq:dual} imply
$\max(H,v+\ell_b)=H+O(\tau_{\mathrm{ref}}+\kappa_G+\eps_Z)$.
When $s_{\mathrm{hyb}}=z_a$, cancel $z_a-s_{\mathrm{hyb}}$ and
discard the other nonpositive terms: this gives at most
$M'-\Delta_H$ plus small losses.
Otherwise $s_{\mathrm{hyb}}\ge v/2$, and
\[
 s_{\mathrm{hyb}}+\ell_b+2e_\lambda/3+(T_d-y)_+/2
 \ge y/4+(T_d-y)_+/2-O(\eps_Z)\ge T_d/4-O(\eps_Z).
\]
The exact saving identity is
\[
\begin{split}
 O/2+\Delta_H+S_0+B_0+T_d/4
 ={}&\frac{2M'+2z_a-1-\ell'+2\Delta_H-\theta_N}{4}\\
 &+\frac{2A_0-N_0+B_0+4S_0}{4}.
\end{split}
\]
The second line is nonnegative. Using $z_a\le\ell'$ gives the
second branch bound
\[
 M'+\frac{1+3\ell'-2M'-2\Delta_H+\theta_N}{4}
       +O(\eps_{\mathrm{ref}}+\kappa_G+\eps_Z).
\]
Here $1+3\ell'-2M'=5\ell-1+d_J$.
Both branches decrease with $\Delta_H$. Extract the kernel saving
from the common profile before squaring or positive enlargement.
Fourier Minkowski, the annular sum, and the finite sector/dyadic sums
therefore give
\begin{equation}\label{eq:Benergy}
 \sum_{0<q_m\ll Q}|B_m^J|^2
 \ll Z^{M'+\pos{(5\ell-1+d_J)/4}+\eps}.
\end{equation}
The powerful-part count is already included here. It must not be
multiplied in a second time. Empty surviving slots lie in the first
branch, apart from vanishing bounded logarithmic offsets.

Throughout $1/6\le\ell\le1/5$ we have
\[
 M'\ge2/5,\quad l_x-d_J\ge11/80,\quad
 l_y-d_J-11b/6\ge1/30.
\]
Thus $Q,Y'\ge1$ eventually and are in fixed polynomial ranges.
Also $P_a=Y'^2/Q=q_{b_*}^{-1}Z^b\ge1$ and
$P_a^2/Y'\ll P_a^{1/6}$. The generic additive Gram result yields
\[
 \sum_{q_m\ll Q}|A_m(Y')|^2
 \ll (1+|\nu|)^{J_{\mathrm{Gr}}}(Q/Y')P_a^{1/6}Z^\eps.
\]
The original Mellin test absorbs this height power.
Cauchy--Schwarz on the same ball, including its $Q^{-1/2}$
normalization, bounds a fixed rescaled tuple by
\[
 (X')^{1/2}P_a^{1/12}
 Z^{\pos{5\ell-1+d_J}/8+\eps}.
\]
There are at most $Z^{d_J+\eps}$ such tuples; their outer coefficients
are $O(Z^{-3d_J/2})$, and shortening $X'$ contributes $Z^{-d_J/2}$.
The total extra exponent is
\[
 -d_J+\frac{\pos{5\ell-1+d_J}}8\le-\frac{7d_J}{8}\le0,
\]
because $5\ell-1\le0$. This pays, rather than deletes, the positive
loss in \eqref{eq:Benergy}. Sum the finitely many subsets to obtain
$X^{1/2}Z^{b/12+\eps}=Z^{(1-\ell)/4-b/6+\eps}$.
Choose the reflection, separation and annular losses within the
prescribed $\eps$, then fix seminorm orders, and finally raise the
threshold. All logarithmic multiplicities are absorbed at this last step.
\end{proof}
For \eqref{eq:geometry} the low exponent equals $C(\sig)$ exactly.

\section{Euler corrections, contours, and the principal signal}
\label{sec:euler}
Let $V,W,D,P_p,P_p^*,H_p$ be the source local factors; $W$ here is a
local Euler factor, distinct from the smooth slot windows. For a
sixth-power-free physical row $u$, write $\chi_\bullet(u)$ for its
source character presentation. More explicitly, with $Q=q_p$ and
$j=v_p(u)\in\{0,\ldots,5\}$, put
\[
\begin{gathered}
 \rho=\chi_p(u/p^j),\quad
 a_p=\overline{\alpha(p)}^3\eta(p)^3,\quad
 b_p=\overline{\alpha(p)}^2\eta(p)^2\overline{\chi_p(4)},\\
 V=Q^{-6z},\quad R=a_p^2Q^{4-6s-6z},\quad
 W_{\mathrm{loc}}=\rho Q^{-w},\quad
 W=\chi_p(u)Q^{-w},\quad D=\eta(p)\overline{\chi_p(u)}Q^{-s}.
\end{gathered}
\]
The source phases have modulus one and $b_p^3=a_p^2$.
Its coefficientwise local table gives
\[
 \begin{split}
 P_p^*&=\frac1{1-R}
 \left\{\frac{R(1-Q^{-1})-\eta(p)(Q-1)Q^{-s-w}V^{1_{j\le1}}}
 {1-V}+J_j\right\},\\
 P_p&=(1-V)^{-1}+1_{j=0}\frac{W_{\mathrm{loc}}}{1-W_{\mathrm{loc}}}
                 +P_p^*,\qquad
 H_p=P_p\frac{(1-V)(1-W)}{1-D},
 \end{split}
\]
where
\[
\begin{array}{c|l}
j&J_j\\ \hline
0&-\eta(p)\rho^{-1}Q^{-s}+W_{\mathrm{loc}}R\\
1&\eta(p)Q^{-s-w}\\
2&a_p\rho^{-3}Q^{3/2-3s}\\
3&-\eta(p)b_p\rho^{-2}Q^{2-3s-w}+b_p^2\rho^{-4}Q^{2-4s}\\
4&-\eta(p)a_p\rho^{-3}Q^{5/2-4s-w}\\
5&-a_p^2Q^{3-6s}.
\end{array}
\]
At $p\mid u$, $D=W=0$, whereas $W_{\mathrm{loc}}$ retains its phase.
The exact
local identity is
\begin{equation}\label{eq:Hlocal}
 H_p-1=
 \frac{D(V+W-VW)-VW+(1-V)(1-W)E_p}{1-D},
 \qquad E_p=P_p^*+D .
\end{equation}
In particular this expression has no $1-W$ denominator.
Use the domains
\[
\begin{split}
\mathcal D_1(\eps_0)=\{&s_r\ge51/100,\ z_r\ge17/50,\
w_r\ge-1/100,\ s_r+w_r\ge1+\eps_0\},\\
\mathcal D_2^*=\{&s_r\ge\sig,\ z_r\ge33/200,\ w_r\ge19/20\},
\end{split}
\]
where subscripts $r$ denote real parts.
The source local table gives, at good primes,
$E_p\ll Q^{4-6s_r-6z_r+\theta}+Q^{1-s_r-w_r-6z_r}$,
where $\theta=(-w_r)_+$ and $Q=q_p$. On $\mathcal D_2^*$,
$\theta=0$.
Substitution in \eqref{eq:Hlocal} gives on $\mathcal D_2^*$
\[
 H_p-1\ll Q^{-1-c_b},\qquad
 c_b=\min(6\sig-401/100,\sig-3/50,47/50)=65199/80000.
\]
At ramified primes it gives
$H_p-1\ll Q^{-c_r}$, where
$c_r=\min(\sig-1/20,3\sig-3/2)=65999/80000$.
On $\mathcal D_1(\eps_0)$ the corresponding exponents are
$-1-\eps_H$ and $-\eps_H$, where
$\eps_H=\min(\eps_0,1/50)>0$.
Strict positive margins yield normal convergence on neighborhoods,
uniform in all imaginary parts, with ramified finite product
$q_u^\eps$.

Define the selected local replacement by
\[
 G_p=
 \frac{(\overline{\eta(p)}Q^s-Q^{-w})(1-V)(1-W)P_p^*
       -Q^{-w}(1-VW)}{1-D}.
\]
Define the full correction by the finite sum
\begin{equation}\label{eq:fulltuple}
 \HH_{\eta,u,Z}=
 \sum_{(p_i)}
 \prod_i\{W_i(q_{p_i}/P_i)q_{p_i}^{z-1}G_{p_i}\}
 \prod_{\substack{p\notin S\\p\notin\{p_1,\ldots,p_K\}}}H_p .
\end{equation}
This definition does not divide by $P_p$ or $H_p$.
The local table bounds $G_p\ll Q^B$ on every fixed real box,
uniformly in heights. Counting the tuples gives
$\HH_{\eta,u,Z}\ll q_u^\eps Z^{B'}$, with $B'$ fixed before any
external integration-by-parts order.

On the absolute starting lines the exact identity
\[
 \overline{\eta(p)}Q^{s+z-1}P_p^*
       -Q^{z-w-1}P_p
 =\frac{1-D}{(1-V)(1-W)}Q^{z-1}G_p
\]
shows that the finite operation of Section~\ref{sec:physical}
extracts the original scalar quotient, once per selected prime.
Consequently it proves, for that independently defined physical sum,
\begin{equation}\label{eq:highidentity}
\begin{split}
 I_{\eta,\mathrm{mod}}(Z)=\frac1{(2\pi i)^3}
 \int_{(3)}\int_{(3)}\int_{(2)}
 &\mathcal W
 \sum_u^{(6)}q_u^{-z}\overline{\xi(u)}\\[-2pt]
 &\frac{\zeta_F^S(6z)L^S(w,\chi_\bullet(u))}
 {L^S(s,\eta\overline{\chi_\bullet(u)})}
 \HH_{\eta,u,Z}(s,w,z)\,dz\,dw\,ds ,
\end{split}
\end{equation}
where all physical row restrictions are retained, and
\[
 \mathcal W=X^{1/2-z}Z^{s+z-1}Y^{w-1}
 e^{(s+z-1)^2}M(z)\widehat W_1(w).
\]
The source coefficientwise identity supplies absolute convergence
on these starting lines. The new local estimates supply its required
analytic domain and all-height majorant.

\subsection{Principal normalization}
Choose $P_0$, before choosing a target, so that
$\exp(C_0\sum_{Q>P_0}Q^{-1-c_b})-1\le1/2$.
Include these small primes in $S$. Enlarging $S$ later only shrinks
this positive majorant. For $u=1$ the unselected product
$H_\eta(s)=\prod_{p\notin S}H_p(s,1,1/6)$ is holomorphic on
$\Rea s>\sig$ and $|H_\eta-1|\le1/2$ there.
Only for this principal row may we now use $G_p/H_p$ throughout
$\mathcal D_2^*$. The four local error powers are
\[
 -s_r,\quad-6z_r,\quad4-5s_r-6z_r,\quad1-w_r-6z_r.
\]
Their maxima on this domain are
$-\sig,-99/100,-21839/16000,-94/100$, respectively. Thus
$G_p/H_p=-1+O(Q^{-\sig})$.

Let $S_i(Z)=\sum_p W_i(q_p/P_i)q_p^{-5/6}$ in the actual ray.
The imported fixed-ray prime asymptotic gives
$S_i\sim c_iP_i^{1/6}/\log P_i$, $c_i>0$.
For our even $K$ set
\[
 A_T(Z)=(-1)^K Z^{-\ell/6}\prod_i S_i(Z).
\]
It is eventually nonzero, $A_T\asymp(\log Z)^{-K}$,
and $A_T^{-1}\ll_\eps Z^\eps$.
For every fixed $0<\mu<\sig\min_i\ell_i$,
\begin{equation}\label{eq:principalfactor}
 \HH_{\eta,1,Z}(s,1,1/6)
 =H_\eta(s)Z^{\ell/6}A_T(Z)(1+R_{\eta,Z}(s)),
 \qquad R_{\eta,Z}(s)\ll Z^{-\mu}.
\end{equation}
The bounds are uniform in heights on the global principal line.
Put
\[
 c_S=\widehat W_1(1)M(1/6)
       \frac{(\operatorname{Res}_{s=1}\zeta_F^S(s))^2}{6}>0.
\]
Moving the principal $w,z$ lines to $1+e,1/6+e$, then
$w=19/20$ across $1$ and $z=33/200$ across $1/6$, stays in
$\mathcal D_2^*$. Keep $s=\bet+e$ throughout these moves.
The reciprocal is globally controlled there; it crosses no zero.
Only the main double residue is moved right to $s=2$.
Define
\begin{equation}\label{eq:f}
 f_\eta(Z)=\frac1{2\pi i}\int_{(2)}
 Z^{C(s)}e^{(s-5/6)^2}\frac{H_\eta(s)}{\LL(s,\eta)}\,ds.
\end{equation}
If $P_\eta$ denotes the full principal row, the residual integrals,
including full two- and three-dimensional tails, satisfy
\begin{equation}\label{eq:principalerrors}
\begin{split}
 \left|P_\eta/(c_S A_T)-f_\eta\right|\ll {}&
 Z^{C(\bet)+(1+h)e-m_w+\eps}
 +Z^{C(\bet)+e-m_z+\eps}\\
 &+Z^{C(\bet)+e-\mu+\eps},\\
 m_w={57497}/{2400000},\qquad
 m_z={32503}/{24000000}.
\end{split}
\end{equation}

\subsection{Whole buffered bins and the outer rows}
\label{sec:outer}
Define
\[
 \bet=\sup\left(\{1/2\}\cup
 \{\Rea\rho:\ L_F(\rho,\eta)=0,\ 1/2\le\Rea\rho\le1,\
                   \eta\ \text{primitive of finite order}\}\right).
\]
Poles are excluded from this set. Assume for contradiction
$\sig<\bet\le7/8$, with the upper bound supplied by $\mathcal R$.
Use the original buffered-bin partition before moving
any row contour. Its floor is $a=51/100$; elsewhere $a\le\bet$.
At $z_0=17/50$ move the whole bin first to
$s=\bet+20e$, $w=1-a-6e$. Remove the outer height box on this
global line, then move the retained bin to $s=a+16e$.
We have $s+w\ge1+10e$ and $s>a+6e$, so $\mathcal D_1(10e)$
and the buffered reciprocal bound apply. The floor need not have
a witness. Any pointwise amplitude or error partition is made
after this whole-bin move. Discarded pieces and joins are
$O_{\eta,N}(Z^{B_\eta}T_1^{-N})$ with $B_\eta$ fixed before $N$.

Here are independent controls for the outer rows, needed to cover
every physical row. Put $d_{\min}=1/100$. On the small-row lines
$(s,w,z)=(\bet+e,1/2,z_0)$ use $\mathcal D_1(1/3)$:
$\sig+1/2>1+1/3$. The old domain $\mathcal D_1(3/8)$ would be
insufficient. The full local table gives $G_p=O(1)$ off the row
and $G_p\ll Q^{1/2}$ on it. Thus each positive selected-slot sum
is at most $q_u^\eps P_i^{z_0}$, without dividing by $H_p$.
Hecke growth contributes $U^{3/5}$, row count $U$, and $q_u^{-z_0}$
contributes $U^{-z_0}$, giving $U^{63/50}$.
The exponent relative to $C(\bet)$ for $U\le Z^{d_{\min}}$ is at most
\begin{equation}\label{eq:small}
 h\frac{13}{75}-\frac{l_y}{2}+\frac{63}{50}d_{\min}
 =-\frac{172249}{2000000}.
\end{equation}
For $U>Z^{h+\zeta}$ keep the full tuple and use absolute lines
$(s,w,z)=(2,2,z_\infty)$, $z_\infty>2$. The bound is
\begin{equation}\label{eq:large}
 Z^{B_0+(h+\zeta)(1+\eps)-\zeta z_\infty+\eps},
 \qquad B_0=\frac{132497}{80000}.
\end{equation}
For any fixed $\zeta>0$, choose $z_\infty$ sufficiently large
before the target. It gives any prescribed fixed saving. These
estimates retain all bounded units and the full absolute row tails.

\section{Actual detector capacities and whole-slot counts}
\label{sec:count}
The plain moment input requires $3/4\le\kappa\le1$ and, if
$\kappa<1$, $\bet\le(1+\kappa)/2$. Under the contradiction above
we may fix $\kappa=3/4$. In general the admitted value would be
$\max(3/4,2\bet-1)$, with excess
$\max(\bet-7/8,0)$. Here that excess is zero.

Write $\delta=2a-1$. For positive amplitude bins use the source bounds
$|Q_i|\le P_i^{g_i+\vartheta}$ and, only if $g_i>0$,
$|Q_i|\ge P_i^{g_i}$, with a prescribed small bin loss
$\vartheta>0$.
Set $g_i=0$ for error slots and for main slots with $|Q_i|\le1$;
the latter need only an upper bound and request no positive slot.
Absorb the preliminary bin losses separately, with
$0\le g_i\le\delta/2$, and let
$q=\sum_i\ell_i g_i/\ell\in[0,\delta/2]$ be their length-weighted
mean; put $x=q/\delta\in[0,1/2]$.
Thus a squared slot product has gain exponent $2q$.
The inverse amplification is on the actual injective map
$(u,a')\mapsto u(a')^6$, with conductor $\ll U$ and width one.
It gives $e(r)=\max(1,(1+5r)/6)$, hence the long count
\begin{equation}\label{eq:longcount}
 L(t)=1-\delta+(5/6-\delta)(t-1).
\end{equation}
For an actual detector witness the imported dyadic construction gives
\[
 r+m\ge t-o(1),\quad r\le t+o(1),\quad
 t-1/2-O(\eps)\le r,\quad 0\le m\le1/2+O(\eps),
\]
with $|M_r|^2\gg U^{\delta r-\eps}$ and
$|S_m|^2\gg U^{\delta m-\eps}$.
The two occurrences of $S_m$ in the plain argument give $|S_m|^4$.
Its slot capacity is $(1-2m)/(6\kappa)$; inverse amplification
has capacity $(1-r)/2$.

In the inverse branch $r\ge r_*$, for a positive requested slot
length $z=(1-r)/2-\nu_0$, decreased by a fixed $\nu_0>0$,
the inverse capacity inequalities include
\[
 1-r-2z\ge2\nu_0,\qquad
 3-2r-8z\ge8\nu_0+9/37-O(\eps).
\]
In the plain branch $r\le r_*$, use
$z=(1-2m)/(6\kappa)-\nu_0$; its inequality is
$1-2m-6\kappa z\ge6\kappa\nu_0$.
The ideal short exponents at $\kappa=3/4$ are
\[
 A_I(r)=1-\delta[x+(1-x)r],\qquad
 S_t(r)=1-\delta[4x/9+(2-8x/9)(t-r)].
\]
Set
\begin{equation}\label{eq:DP}
 D_x=3-17x/9,\quad P_x=(2-8x/9)(1-x),\quad
 J_x=(5/6-\delta)D_x+\delta P_x.
\end{equation}
Their crossing is
$r_*=[(2-8x/9)t-5x/9]/D_x$ and gives
\[
 R_{\mathrm{short}}=1-\delta+
          \delta(P_x/D_x)(3/2-t).
\]
On $1\le t\le3/2$, $r_*\ge23/37$,
$1/3\le t-r_*\le1/2$, and $r_*(3/2)=1$.
The maximal inverse and plain capacities are respectively
$7/37$ and $2/27$. Near saturated $r\ge1$ use amplification;
near $m=1/2$ use the zero-slot plain moment. A capacity at most
$\nu_0$ costs at most $2\delta\nu_0$ or $6\delta\nu_0$,
respectively, in the corresponding count.

The actual normalized prime factors are
\begin{equation}\label{eq:Qi}
 Q_i(u;z)=P_i^{-1/2}
  \sum_{p\in\mathcal P_i(Z)}
  \overline{\chi_p(u)}W_i(q_p/P_i)(q_p/P_i)^{z-1}.
\end{equation}
Primes dividing $u$ have zero character value. The fixed presentation
is $\psi=\nu\overline{\chi_\bullet(u)}$, with
$\nu\in\Theta=\langle\eta,\widehat T\rangle$, a fixed finite group
for the target. Its row-independent prime coefficients are
$\overline{\nu(p)}1_T=|T|^{-1}\sum_{\theta\in\widehat T}
(\overline{\nu}\theta)(p)$, a finite combination of members of $\Theta$.
In a plain moment, conjugate both witnesses
and all selected factors; the coefficients become $\nu(p)1_T$.
The supports and masks are unchanged.

The source inducing-character exclusion also remains valid:
at a prime with sixth-free ramification label $j=1,\ldots,5$,
the local order $6/\gcd(6,j)>1$ cannot be cancelled by an unramified
fixed presentation. The remaining $S$-supported units form a finite
set, eventually outside selected rows $U\ge Z^{1/2}$; outer estimates
retain them.

For $U=Z^d$ the whole-slot lengths are $w_i=\ell_i/d$.
Greedy selection of full positive slots loses at most one fraction,
so its gain satisfies
$|\prod Q_i|^2\ge U^{2qz-\delta\max_iw_i}$.
If the positive supply is insufficient, taking it all still gives
gain at least $qz$. For the selected range $d\ge1/2$,
$w_i\le2\ell_i$. There is strict supply:
\[
 \frac{\ell}{h}-\frac7{37}=\frac{57659}{3607833}>0,\qquad
 0<\zeta<5\ell-h=\frac{2521}{120000}
 \ \Longrightarrow\ \frac{\ell}{h+\zeta}>\frac15>\frac7{37}.
\]
The latter gap is $2/185$. Choose the uniform plain-moment mesh
before $K$, and then fix even $K$ with
$2\ell/K<\min(\text{mesh},\text{rounding allowance},1/185)$.
The actual $w_i$ are smaller than this mesh, and every physical
slot norm is at most $2U^{\mathrm{mesh}}$.
The plain moment's internal pool starts at
$U^{\mathrm{allowed\ width}/3}/2$.
Choose $\mathrm{mesh}<\mathrm{allowed\ width}/6$, giving a strict
exponent gap and eventual disjointness from that pool.
Annular offsets are $O_K(1/\log U)$
and enter only after all strict margins have been fixed.

The admitted marked, inverse and plain moments now count the
actual rows, with the original witnesses and whole selected slots:
\[
 \#\mathcal B\ll
 U^{\max(R_{\mathrm{short}},L(t))+\eps}(1+T_1)^A.
\]
Take
\begin{equation}\label{eq:R}
 t=1+\frac{\delta P_x}{2J_x},\qquad
 R_*=1-\delta+
      \frac{(5/6-\delta)\delta P_x}{2J_x}.
\end{equation}
Here $35/54\le J_x\le5/2$, $D_x\ge37/18$, $P_x\ge7/9$,
and $1<t<3/2$. This balances the two counts and proves
$\#\mathcal B\ll U^{R_*+\eps}(1+T_1)^A$.
The no-slot choice $t=1$ instead gives exponent $1-2\delta/3$.
The floor uses the trivial exponent $1$, without an invented witness.

\section{Joint errors and a continuous high-side saving}
\label{sec:high}
On the retained dynamic line
$(s,w,z)=(a+16e,1-a-6e,z_0)$,
the local product satisfies good/ramified bounds
$Q^{-1-10e}$ and $Q^{-10e}$.
After fixing $e$, enlarge the target-independent small-prime cutoff.
Only on this dynamic line may we factor by $H_p$ and expand
\[
 \HH=\mathcal H_{\eta,u}
       \sum_I\prod_{i\in I}D_i\prod_{i\notin I}
                       (-P_i^{z-1/2}Q_i),
\]
where $D_i$ denotes the local error sum. The global contour
argument used the full expression \eqref{eq:fulltuple}.
Off the row its four error exponents are
$-a-4e,-2.04+12e,1.96-5a-68e,a-2.04+12e$,
all at most $-0.51$ for sufficiently small $e$.
At a ramified prime the normalized table has, at $e=0$, powers
\[
 a-2,\quad 1/2-2a,\quad (-a,1-3a),\quad
 1/2-2a,\quad 2-5a ,
\]
with $e$-coefficients $6,-32,(-26,-48),-42,-80$.
These are strictly less than $-1/2$, except for the strict
$j\ge2$ term whose bound is $Q^{-w_r}$.

All such strict terms must share one conductor budget.
For the entire triangle-expanded set of distinct strict primes,
the primitive conductor satisfies
$q_{\mathfrak f_u}\ll_S q_u\prod_p Q_p^{-(j_p-1)}$.
Apply the numerator functional equation once, with
$A^*=1/2-w_r=a-1/2+6e$. Its reflected primitive value is at
$a+6e$ and is covered by the buffered reciprocal control,
including the conjugate presentation. Deleted factors cost
$U^{6e+\eps}$. Each strict prime then contributes at most
\[
 Q_p^{z_0-w_r-(j_p-1)A^*}\le Q_p^{z_0-1/2},
 \qquad j_p\ge2.
\]
Restoring its deleted strict factor costs $1$.
Consequently, for every error subset $I$, including the empty one,
\[
 \left|L^S(w,\chi_\bullet(u))\prod_{i\in I}D_i\right|
 \ll U^{\delta/2+O(e)+\eps}(1+T_1)^A
       \prod_{i\in I}P_i^{z_0-1/2+O(e)+\eps}.
\]
Assign gain $g_i=0$ to errors and $0\le g_i\le\delta/2$ to
main factors, and let $q=\sum_i\ell_i g_i/\ell$.
The joint pointwise exponent is therefore
$U^{\delta/2+\mathrm{loss}}Z^{\ell(z_0-1/2)+q\ell+\mathrm{loss}}$.
The count of Section~\ref{sec:count} uses this same $q$ over
all slots. No second conductor saving is charged.

After the source contour accounting, the exponent relative to
$C(\sigma)$, before small losses and height powers, is
\begin{equation}\label{eq:E}
 E_\sigma(d)=K_\sigma+(1/2+\ell)\delta+\ell q
       -h(1-R)+(d-h)(R+\delta/2-z_0),
 \qquad K_\sigma=2/3+\ell+b/6-\sigma .
\end{equation}
Equivalently the unfactored accounting is
$a-\sigma+h(z_0-1/6)-a l_y-\ell/2+q\ell+
d(R+\delta/2-z_0)$.
The contour's explicit additional loss is
$(16-6l_y)e+(1+d)\eps_c+d\eps_d+\eps_p$,
besides previously fixed count, normalization and height losses.

\begin{lemma}[Continuous exact certificate]\label{lem:certificate}
At the geometry \eqref{eq:geometry}, with $R=R_*$ and
$q=\delta x$, $0\le x\le1/2$, $1/50\le\delta\le3/4$,
\[
 E_{\sig}(h)\le-\frac{307}{11016000}=:-m_{\mathrm{ad}}.
\]
\end{lemma}
\begin{proof}
First set $\ell_0=1/6$, $\sigma_0=7/8$, $h_0=13/16$,
and denote the endpoint exponent by $E_0$. Put
$y=1/2-x$, $v=51+41y$, and use $J_x$ of \eqref{eq:DP}.
Direct expansion gives the polynomial identity
\begin{align*}
 10368vJ_x(-E_0)
 ={}&(3+5y)\{(4v\delta-79)^2+49\}\\
 &+4y\{4v\delta[(1+3y)(15+32y)\delta+9-13y]
                       +265+3485y\}.
\end{align*}
Every term is nonnegative for $0\le y\le1/2$ and
$0\le\delta\le3/4$, since $9-13y\ge5/2$.
Also $v\le17(3+5y)$ and $J_x\le5/2$, hence
$-E_0\ge49/440640$.
The inequality $3P_x\le2D_x$ implies $R_*\le1-2\delta/3$.
Along $\sigma(\ell)=11/12-\ell/4$ one obtains exactly
\[
 E_{\sigma(\ell_0+t)}(h_0+3t/2)-E_0
 =t(-1/4+\delta+q+3R_*/2)\le5t/3.
\]
For $t=1/20000$ subtraction yields
$49/440640-1/12000=307/11016000$.
This already includes the change of reference exponent $C(\sigma)$;
no further deduction of $1/80000$ is appropriate.
\end{proof}

For $1/2\le d\le h$, the slope in \eqref{eq:E} is positive,
so the endpoint controls the whole interval.
The floor, with $\delta=1/50$, $R=1$, $q\le1/100$,
has $E_{\sig}(h)\le-4327/750000$.
For $d_{\min}\le d\le1/2$ the common no-slot bound
$R=76/75-2\delta/3$ covers the floor too. Its $d$ slope is
$101/150-\delta/6>0$; the geometry variation at fixed $d$
has derivative at most $13/50+\delta/4+q\le177/200$.
Direct substitution gives the uniform bound
$E_{\sig}(d)\le-120907/36000000$.
Choose
\begin{equation}\label{eq:zeta}
 \zeta=m_{\mathrm{ad}}/32=307/352512000.
\end{equation}
It satisfies the supply restriction in Section~\ref{sec:count}.
The slope is at most $2$, so extending to $h+\zeta$ costs at
most $m_{\mathrm{ad}}/16$ and leaves
$307/11750400$.
Relative to $C(\bet)$ all these exponents decrease further by
$\Delta_1=\bet-\sig$.
Together with \eqref{eq:small} and \eqref{eq:large},
these estimates cover every nonprincipal physical dyad.

\section{Parameter order, height closure, and continuation}
\label{sec:closure}
\begin{proof}[Proof of Theorem~\ref{thm:main}]
Under the contradiction $\sig<\bet\le7/8$, fix
$\Delta_1=\bet-\sig>0$.
All real exponents below are fixed before the target character.
First prescribe the low and inverse-normalizer losses so their sum
is at most $\Delta_1/2$. Proposition~\ref{prop:low} then gives, for
\[
 J_\eta(Z)=I_{\eta,\mathrm{mod}}(Z)/(c_S A_T(Z)),
 \qquad \omega=\Delta_1/2,
\]
the bound $|J_\eta(Z)|\ll_\eta Z^{C(\sig)+\omega}$.

For the high side first choose moment losses, strict capacity
decrements, the uniform plain mesh and rounding allowances within
small fixed portions of $m_{\mathrm{ad}}$. Next choose the fixed
even $K$ as in Section~\ref{sec:count}, then the windows. Only now put
$\mu=\sig\ell/(2K)>0$. Choose amplitude widths and preliminary
prime losses after $K$, and choose $e>0$ respecting the detector's
bounded presaturation ranges and all local domains.
The total central real, dyadic, count and normalizer losses are
less than $m_{\mathrm{ad}}/4$.
Make the principal losses satisfy
\[
 (1+h)e+\eps_{\mathrm{pr}}<m_w/4,\quad
 e+\eps_{\mathrm{pr}}<m_z/4,\quad
 e+\eps_{\mathrm{pr}}<\mu/4.
\]
Fix \eqref{eq:zeta}, then $z_\infty$ in \eqref{eq:large};
choose the small-row losses below half the margin in \eqref{eq:small}.
The remaining central saving is at least $m_{\mathrm{ad}}/2$.
Let
\[
 m_0=\min(m_{\mathrm{ad}},m_w,m_z,\mu,172249/2000000),
 \qquad m=m_0/4>0.
\]
All ranges now retain at least $m_0/2$. The earlier mesh was paid
by $m_{\mathrm{ad}}$, and is not reselected using the smaller,
$K$-dependent $\mu$. This avoids a circular order of choices.

Choose a target $\eta$, fix its final arithmetic data and $S$,
and then fix all internal moment, Sobolev and seminorm orders.
The same $S$ is used in every physical sum, identity, calibration,
$c_S$, $A_T$, $H_\eta$ and $f_\eta$.
The preceding estimates give finite exponents $A_\eta,B_\eta$,
independent of any later external tail order $N$, such that
\begin{equation}\label{eq:highestimate}
 |J_\eta-f_\eta|
 \ll_{\eta,N} Z^{C(\bet)-m}(1+T_1)^{A_\eta}
                  +Z^{B_\eta}T_1^{-N}.
\end{equation}
The physical $\Ima z$, detector witnesses and prime frequencies use
the original fixed matrix's common cumulative $T_1/2$ allocation.
They do not get separate newly enlarged buffers.

Let $\eps_{\mathrm{ht}}>0$ be the minimum of the finitely many
specified detector height allowances, and let
$A_{\mathrm{ht},\eta}\le A_\eta$ dominate their powers. Set
\[
 \tau_{0,\eta}=
 \frac{d_{\min}\eps_{\mathrm{ht}}}{20(1+A_{\mathrm{ht},\eta})}.
\]
For $T_1=Z^\tau$ with
$0<\tau\le\min(d_{\min}/100,\tau_{0,\eta})$ and
$U\ge Z^{d_{\min}}$, we have eventually
$T_1\le U^{1/100}$ and
$(1+T_1)^{A_{\mathrm{ht},\eta}}\le U^{\eps_{\mathrm{ht}}/10}$.
All literal detector height hypotheses thus hold; their crude
$U^{-189/100+o(1)}T_1^2$ error tends to zero.
Small and absolute large rows do not need this buffered cutoff.
Increasing $N$ in \eqref{eq:highestimate} changes only external
test seminorms, not internal orders or the real powers.

Now choose
$0<\tau_\eta\le\min(d_{\min}/100,\tau_{0,\eta},
m/[4(A_\eta+1)])$, then an integer $N_\eta$ with
$B_\eta-N_\eta\tau_\eta<C(\bet)-m/2$, and finally the threshold.
It follows that
\[
 |J_\eta-f_\eta|\ll_\eta Z^{C(\bet)-\sigma_{\mathrm{hi}}},
 \qquad \sigma_{\mathrm{hi}}=m/2.
\]
Although $\tau_\eta,N_\eta$ and constants depend on the target,
$\omega,\sigma_{\mathrm{hi}}$ do not. Neither $J_\eta$ nor
$f_\eta$ contains $T_1$ in its definition.
With $\eps_*=\min(\Delta_1/2,\sigma_{\mathrm{hi}})>0$, this gives
$f_\eta(Z)\ll_\eta Z^{C(\bet)-\eps_*}$ at infinity.
At zero, move \eqref{eq:f} right to any fixed $B>2$.
The Gaussian and the absolutely convergent reciprocal Euler product
give $f_\eta(Z)\ll_{\eta,B} Z^{B-11/16}$ for $0<Z\le1$.
Consequently
\[
 F_\eta(s)=\int_0^\infty f_\eta(Z)Z^{-C(s)}\,\frac{dZ}{Z}
\]
is locally uniformly convergent and holomorphic for
$\Rea s>\bet-\eps_*$.
Writing $Z=e^u$, Fourier inversion on $\Rea s=2$ identifies it with
$e^{(s-5/6)^2}H_\eta(s)/\LL(s,\eta)$.
The established integrability justifies the inversion.
The identity theorem extends this identity to $\Rea s>1$.

Since $\bet-\eps_*>\sig$ and $|H_\eta|\ge1/2$,
$e^{-(s-5/6)^2}F_\eta/H_\eta$ is a holomorphic continuation of
$1/\LL(s,\eta)$ to $\Rea s>\bet-\eps_*$.
The definition of the supremum supplies some target with a zero
whose real part exceeds $\bet-\eps_*$; attainment of $\bet$ is
unnecessary. Its reciprocal has a pole, a contradiction.
The deleted finite Euler factors are nonzero for $\Rea s>0$.
Thus $\bet\le\sig$, proving the primitive Hecke assertion and,
by restoring finite Euler factors, the imprimitive assertion.

Finally the imported quadratic transfer gives
\[
 L_F^S(s,\chi\circ\N)
     =L^S(s,\chi)L^S(s,\chi\chi_{-3}).
\]
In $0<\Rea s<1$ both Dirichlet factors are holomorphic, so their
zeros cannot be cancelled by poles. Above $1$ their Euler products
are nonzero. At $1+it$, $t\ne0$, the same product excludes zeros.
At $1$ the only possible pole-zero pairing is between the principal
character and $\chi_{-3}$, but
$L(1,\chi_{-3})=\pi/(3\sqrt3)>0$.
This proves the stated Dirichlet conclusion, including zeta.
\end{proof}

\clearpage
\appendix
\section{Verification and reproducibility}
\label{app:verification}
The derivation was developed in repository notes 441--445
\cite{Project}. The exact audit
\src{scripts/hybrid_high_continuation_exact_audit.py}
uses rational arithmetic for the geometry, Euler margins,
capacity supply, all physical row ranges and the continuous
completed-square identity. Its output is
\src{output/hybrid-high-continuation-exact-audit.json}.
The polynomial identity in Lemma~\ref{lem:certificate}, rather than
a grid sample, proves the continuous central inequality.
The audit also checks 16,728 finite frequency-allocation cases.
These calculations check the stated algebra and bookkeeping;
they do not certify the imported analytic theorems.

The pinned source canonical SHA-256 is
\begin{center}\small\ttfamily
42a5ee0febca59fd1def55cfd6c6808c322ef\\
4237d7726303ea52f711deac6a3 .
\end{center}
The research checkpoint for notes 441--445 is
\src{96fcb27a0b8788ad8a346238843906bce0285479}.
Independent source-reading reviews and the final paper build
manifest are stored in \src{reviews/2026-10-07/}.
The manifest binds the final source, PDF and review hashes.

No external Lean elaboration, kernel build, or Comparator run was
performed for this improvement. The source's generic results and
its whole Hecke seven-eighths theorem are explicit imported inputs.
Independent verification of that package and a formalization of
the present extension remain separate tasks. This paper confines
its mathematical assertion to Theorem~\ref{thm:main}.

\begin{thebibliography}{9}
\bibitem{OAI}
OpenAI, \emph{The Quasi-Riemann Hypothesis},
September 30, 2026 version, source at commit
\src{adc7f1241b42e322a6451854ab7e4b4c146bf78a},
\href{https://github.com/openai/math/blob/adc7f1241b42e322a6451854ab7e4b4c146bf78a/preprints/The-Quasi-Riemann-Hypothesis-September-30-2026/build/paper.tex}{pinned source}.
\bibitem{Project}
RH--Weil research project, \emph{Variable-slot low bound, shared
contours, actual detector capacities, joint errors, and family
continuation}, notes 441--445 and their independent reviews,
October 7, 2026, checkpoint
\href{https://github.com/lixiang90/RH-Weil/tree/96fcb27a0b8788ad8a346238843906bce0285479}{96fcb27a0b8788ad8a346238843906bce0285479}.
\end{thebibliography}
\end{document}
